\clearpage
\SourceLeaf{279}
\renewcommand{\BookRunningTitle}{Klassenanzahl und Determinante}
\renewcommand{\Article}[1]{\par\medskip\Needspace{4\baselineskip}\pdfbookmark[1]{Art. #1}{nexus-prior-art-#1}\centerline{\textbf{#1.}}\nobreak\smallskip}
\pdfbookmark[0]{Klassenanzahl und Determinante. I}{nexus-prior}
\begin{center}
[I.]\par\medskip
\Large\bfseries ÜBER DEN ZUSAMMENHANG ZWISCHEN DER ANZAHL DER KLASSEN, IN DIE BINÄRE QUADRATISCHE FORMEN ZERFALLEN, UND IHRER DETERMINANTE.\par\medskip
\normalsize ERSTE ABHANDLUNG\par\medskip
DER KÖNIGLICHEN GESELLSCHAFT VORGELEGT \mbox{1834\ .\ .\ .\ .}
\end{center}
\BookRule{.18\textwidth}

\Article{1}
Dreiunddreißig Jahre sind nun vergangen, seit wir die Prinzipien des bemerkenswerten Zusammenhangs entdeckten, dem diese Abhandlung gewidmet ist, wie schon am Schluss der \emph{Untersuchungen über höhere Arithmetik} angekündigt wurde. Andere Beschäftigungen hatten uns jedoch lange von dieser Untersuchung abgehalten, bis wir in jüngerer Zeit zu ihr zurückkehren und sie durch neue Bemühungen erweitern konnten. Da aber dieser neue Teil der höheren Arithmetik den Umfang einer einzigen Abhandlung übersteigt, soll diese erste den Formen mit negativen Determinanten gewidmet sein; die Formen mit positiven Determinanten, die eine ganz eigene Behandlung erfordern, müssen einer anderen Abhandlung vorbehalten bleiben.

\Article{2}
Die Grundlage des ganzen Gegenstands bildet eine besondere Untersuchung über die Anzahl aller Kombinationen ganzzahliger Werte, die zwei unbestimmte ganze Zahlen $x,y$ innerhalb eines vorgeschriebenen Bereichs annehmen. Offenbar lässt sich diese Aufgabe auch geometrisch darstellen, indem man die Anzahl der \emph{komplexen Zahlen} ermittelt, deren Darstellung in eine vorgeschriebene Figur fällt. Die Beschaffenheit der Figur hängt von der Beschaffenheit der sie umgebenden Linie ab und somit entweder von einer einzigen Gleichung zwischen den Koordinaten $x,y$ (wenn der Rand eine in sich zurückkehrende Kurve ist) oder von mehreren solchen
\SourceLeaf{280}
Gleichungen (wenn er .aus mehreren gekrümmten oder geraden Teilen besteht); dabei bleibt es unserer Wahl überlassen, ob wir Punkte, die komplexen ganzen Zahlen entsprechen und etwa auf dem Rand selbst liegen, mitzählen oder ausschließen wollen.

\Needspace{5\baselineskip}
Bei der analytischen Darstellung der Aufgabe lassen sich diese Begrenzungsbedingungen stets so ausdrücken, dass eine oder mehrere gegebene Funktionen der Variablen $x,y$, nämlich $P,Q,R$ usw., positive oder nichtnegative Werte annehmen müssen (je nachdem, ob der Wert $0$ ausgeschlossen oder zugelassen wird).

Ist beispielsweise die vorgeschriebene Figur ein Kreis mit dem Radius \mbox{$=\surd A$}, und liegt sein Mittelpunkt in einem Punkt, der einer komplexen ganzen Zahl entspricht, so lautet die analytische Bedingung, dass $A-xx-yy$ nicht negativ ist, sofern wir, was wir immer voraussetzen werden, die auf dem Rand selbst liegenden Punkte beibehalten wollen. Ist die Figur ein Dreieck, so müssen die drei linearen Funktionen $ax+by+c$, $a'x+b'y+c'$, $a''x+b''y+c''$ nichtnegative Werte haben, und ebenso in anderen Fällen.

\Article{3}
Eine \emph{exakte} Lösung der Aufgabe muss im Allgemeinen so vorgehen, dass zunächst aus der Beschaffenheit der Bedingungen die eine Variable, beispielsweise $x$, auf Grenzen beschränkt wird, zwischen denen die einzelnen ganzzahligen Werte nacheinander durchlaufen werden. Es ist zu ermitteln, wie viele ganzzahlige Werte der anderen Variablen $y$ jedem davon entsprechen; diese Anzahlen sind dann zusammenzuzählen. In besonderen Fällen wird es meist besondere Verfahren geben, um die Arbeit abzukürzen.

Ist die Figur beispielsweise wie oben ein Kreis mit dem Radius \mbox{$=\surd A$}, so sei $r$ die nächstkleinere ganze Zahl unterhalb von $\surd A$ oder $\surd A$ selbst, falls $A$ ein Quadrat ist. Ebenso seien $r',r'',r'''$ usw. $r^{(r)}$ die nächstkleineren ganzen Zahlen unterhalb von $\surd(A-1)$, $\surd(A-4)$, $\surd(A-9)$ usw. bis zu $\surd(A-rr)$. Dann ist die gesuchte Anzahl
\[
\begin{aligned}
&=2r+1+2(2r'+1)+2(2r''+1)+2(2r'''+1)+\text{usw.}\\
&=1+4r+4r'+4r''+4r'''+\text{usw.}+4r^{(r)}
\end{aligned}
\]

In diesem Beispiel ist die folgende Methode kürzer. Es sei $q$ die nächstkleinere ganze Zahl unterhalb von $\surd\tfrac12 A$ (oder dieser gleich, falls sie ganzzahlig ist); ferner seien $r^{(q+1)}$, $r^{(q+2)}$, $r^{(q+3)}$ usw. die nächstkleineren ganzen Zahlen unterhalb von $\surd(A-(q+1)^2)$, $\surd(A-(q+2)^2)$, $\surd(A-(q+3)^2)$ usw. bis zu $\surd(A-rr)$. Dann ist die gesuchte Anzahl
\[
=4qq+1+4r+8(r^{(q+1)}+r^{(q+2)}+r^{(q+3)}+\text{usw.}+r^{(r)})
\]

\SourceLeaf{281}
Nach dieser Formel wurden die folgenden Anzahlen ermittelt:
\[
\begin{array}{r|r||r|r||r|r}
A&&A&&A&\\\hline
100&317&1000&3149&10000&31417\\
200&633&2000&6293&20000&62845\\
300&949&3000&9425&30000&94237\\
400&1257&4000&12581&40000&125629\\
500&1581&5000&15705&50000&157093\\
600&1885&6000&18853&60000&188453\\
700&2209&7000&21993&70000&219901\\
800&2521&8000&25137&80000&251305\\
900&2821&9000&28269&90000&282697\\
1000&3149&10000&31417&100000&314197
\end{array}
\]

\Article{4}
Für unser Vorhaben ist keine exakte Bestimmung erforderlich, sondern die Suche nach einem Ausdruck, der sich der genauen Anzahl beliebig weit nähern kann, während die Grenzen ins Unendliche erweitert werden. Da dies aber etwas Unbestimmtes enthält, müssen wir die Sache zunächst genauer darlegen.

Wir nehmen also an, dass die Funktionen $P,Q,R$ usw. außer den Variablen $x,y$ ein konstantes Element $k$ enthalten, sodass die einzelnen $P,Q,R$ usw. homogene Funktionen der drei Größen $x,y,k$ sind. Die durch die Gleichungen $P=0,Q=0,R=0$ usw. bestimmte Figur hängt damit von $k$ ab, und zwar entsprechen verschiedenen Werten von $k$ ähnliche und bezüglich des Koordinatenursprungs ähnlich gelegene Figuren; entsprechende lineare Abmessungen sind den Werten von $k$, die Flächeninhalte den Werten von $kk$ proportional. Bezeichnen wir nun die Anzahl der Punkte innerhalb der Figur mit $M$, ihren Flächeninhalt mit $V$, so müssen offenbar $M$ und $V$ mit wachsendem $k$ wachsen; wächst $k$ ins Unendliche, so nähern sich $M$ und $V$ dem Verhältnis der Gleichheit beliebig weit. Oder, wenn du elementare Klarheit verlangst: Ist eine Größe $\lambda$ gegeben, wie klein sie auch sei, so lässt sich stets eine solche Grenze angeben, dass für jeden Wert von $k$, der diese Grenze übersteigt, $\frac MV$ sicher zwischen $1-\lambda$ und $1+\lambda$ liegen muss. Nach der üblichen Ausdrucksweise darf man dies so angeben: Es wird $M=V$ für einen unendlichen \mbox{Wert von $k$}.

In unserem Beispiel gilt die erforderliche Bedingung, wenn wir $k=\surd A$ setzen, und die Kurve wird ein Kreis, dessen Flächeninhalt $=\pi A$, wobei $\pi$ den halben Umfang eines Kreises mit dem Radius $=1$ bezeichnet. Die oben angegebenen Zahlen lassen die Konvergenz deutlich erkennen.

\SourceLeaf{282}
Wenn es übrigens der Mühe wert wäre, könnten wir den Beweis dieses Satzes leicht mit der Strenge der Alten durchführen; an dieser Stelle haben wir ihn jedoch lieber unterdrückt und eilen zu schwierigeren Dingen weiter.

\Article{5}
In dieser Abhandlung wird der Rand durch eine \emph{einzige} Gleichung der Form $axx+2bxy+cyy=A$ ausgedrückt, wobei $a,b,c$ ganze Zahlen sind und $bb-ac$ eine negative Zahl ist, die wir $=-D$ setzen. Offenbar ist die die Figur begrenzende Kurve eine Ellipse, und es ist leicht zu sehen, dass die Quadrate der Halbachsen die Wurzeln der Gleichung
\[
(ac-bb)qq-(a+c)Aq+AA=0\quad\text{oder}\quad
=A\left(\frac{a+c\pm\surd(4bb+(a-c)^2)}{2(ac-bb)}\right)
\]
sind.

Das Produkt dieser Wurzeln ist $\frac{AA}{ac-bb}=\frac{AA}D$, also der Flächeninhalt der Ellipse $=\frac{\pi A}{\surd D}$. Daraus folgt, dass die Anzahl aller Kombinationen ganzzahliger Werte von $x,y$, für die $axx+2bxy+cyy$ den Wert $A$ nicht übersteigt, mit wachsendem $A$ dem Wert $\frac{\pi A}{\surd D}$ immer näher kommt und für unendliches $A$ diesem Wert gleichgesetzt werden muss. Dabei ist offenbar gleichgültig, ob die Kombination $x=0,y=0$ zu den übrigen mitgezählt oder ausgeschlossen wird. Somit ist die gesuchte Anzahl (im letzteren Sinn) nichts anderes als die Summe der Anzahlen der Darstellungen der einzelnen Zahlen $1,2,3\ldots A$ durch die binäre quadratische Form $axx+2bxy+cyy$; und da einige dieser Zahlen durch diese Form überhaupt nicht dargestellt werden können, andere mehr, andere weniger Darstellungen zulassen, ist die Größe $\frac\pi{\surd D}$ als mittlerer Wert der Anzahl der Darstellungen einer unbestimmten positiven Zahl durch eine beliebige Form zu betrachten, deren Determinante $=-D$.
